\subsection{基数的运算}

定义3。设 \((\mathfrak{a}_t)_{t \in I}\) 是一族基数。诸集合 \(\mathfrak{a}_t\) 的积（相应地，和）的基数称为诸基数 \(\mathfrak{a}_t\) 的基数积（相应地，基数和），并且用 \(\underset{i \in I}{\mathsf{P}} \mathfrak{a}_t\) ( \(\text{分别 } \sum_{i \in I} \mathfrak{a}_t\) ) 表示.

只要没有混淆的风险，我们就简称为“积”和“和”，以代替“基数积”和“基数和”，并记作 \(\prod_{i\in I}a_{i}\) 以代替 \(P_{i\in I}a_{i}\) （参见习题2）。

命题4。设 \((\mathbf{E}_t)_{i \in I}\) 为一个集合族，P为它们的积，S为它们的和，并令 \(a_t\) 为 \(\mathbf{E}_t\) 的基数。那么 P（相应地，S）的基数就是族 \((a_t)_{i \in I}\) 的基数积（相应地，基数和）.

因为存在从 P（相应地，S）到诸集合 \((\mathfrak{a}_{t})\) 的积（相应地，和）上的一个双射 (第二章，§4，第8号，命题10，以及§5，第7号，命题11的推论)。

推论. 如果 \((\mathbf{E}_{t})_{t\in\mathbf{I}}\) 是任何集合族，其并集 \(\bigcup_{t\in I}E_{t}\) 的基数至多等于和 \(\sum\) Card \((\mathbf{E}_{t})\) .

因为存在一个从诸 \(E_{t}\) 的和 S 到 \(E_{t}\) 的并集上的映射 (第二章，§4，第8号)；因此，该推论由命题3和命题4得出。

命题5。(a) 设 \((\mathfrak{a}_i)_{i\in I}\) 为一族基数，并令 \(f\) 为集合 \(\mathbf{K}\) 到索引集 \(\mathbf{I}\) 上的一个双射。则

\[
\sum_ {x \in \mathbb {K}} a _ {f (x)} = \sum_ {i \in I} a _ {i}, \quad \underset {x \in \mathbb {K}} {\mathsf {P}} a _ {f (x)} = \underset {i \in I} {\mathsf {P}} a _ {i}.
\]

\begin{enumerate}
\def\labelenumi{(\alph{enumi})}
\setcounter{enumi}{1}
\tightlist
\item
  令 \((\mathfrak{a}_t)_{t\in \mathbf{I}}\) 为一族基数，且设 \((\mathrm{J}_{\lambda})_{\lambda \in \mathbf{L}}\) 为 \(\mathbf{I}\) 的一个划分。则
\end{enumerate}

\[
\sum_ {i \in I} a _ {i} = \sum_ {\lambda \in L} \left(\sum_ {i \in J _ {\lambda}} a _ {i}\right), \quad P _ {i \in I} a _ {i} = P _ {\lambda \in L} \left(P _ {i \in J _ {\lambda}} a _ {i}\right)
\]

（“和与积的结合律”）。

\begin{enumerate}
\def\labelenumi{(\alph{enumi})}
\setcounter{enumi}{2}
\tightlist
\item
  令 \(((\mathfrak{a}_{\lambda t})_{t\in J_{\lambda}}\lambda \in L\) 为一族（由 L 索引的）基数族，且设 \(I = \prod_{\lambda \in I}J_{\lambda}\) 。则
\end{enumerate}

\[
\underset {\lambda \in \mathbf {L}} {\mathbf {P}} \left(\sum_ {i \in \mathrm{J} _ {\lambda}} a _ {\lambda i}\right) \underset {f \in \mathbf {I}} {\sum} = \left(\underset {\lambda \in \mathbf {L}} {\mathbf {P}} a _ {\lambda , f (\lambda)}\right)
\]

（“积对和的分配律”）。

(a)中的关系由集合的并集和乘积的类似公式得出，因为 \(f\) 是双射意味着如果 \((\mathbf{X}_t)_{i \in \mathbf{I}}\) 是一族两两不相交的集合，该族的元素 \((\mathbf{X}_{f(\mathbf{x})})_{\mathbf{x} \in \mathbb{K}}\) 也是两两不相交的（参见第二章，§4，第1号，命题1以及§5，第3号，命题4）。

¶ (b)中的关系是并集与乘积的结合律公式（第二章，§4，第2号，命题2及§5，第5号，命题7）以及交集对并集的分配律（第二章，§5，第6号，命题8）的直接推论，这表明如果 \((\mathbf{X}_{t})_{i\in I}\) 是一族互不相交的集合，则该族的元素

\[
\left(\bigcup_ {i \in J _ {\lambda}} X _ {i}\right) _ {\lambda \in L}
\]

也是互不相交的。

最后，(c) 由乘积对并集和交集的分配性得出（第二章，§5，第6号，命题9及推论1）。

令 \(\mathfrak{a}\) 与 \(\mathfrak{b}\) 为两个基数。如果 \(\mathbf{I}\) 是一个由两个不同元素组成的集合（例如基数2），则存在一个映射 \(f\) （从 \(\mathbf{I}\) 到 \(\{\mathfrak{a}, \mathfrak{b}\}\) 上的映射），这定义了一个基数族。该族的和与积仅依赖于a和b（根据命题5(a)）；这些基数分别称为a与b的和与积，并记为 \(\mathfrak{a} + \mathfrak{b}\) 与 \(\mathfrak{a}\mathfrak{b}\) 。类似地，对于三个或更多基数的和与积也是如此。于是命题5蕴含以下推论：

推论。设a、b、c为基数。则

\begin{enumerate}
\def\labelenumi{(\arabic{enumi})}
\tightlist
\item
\end{enumerate}

\[
\mathfrak {a} + \mathfrak {b} = \mathfrak {b} + \mathfrak {a}, \quad \mathfrak {a}\mathfrak {b} = \mathfrak {b}\mathfrak {a};\tag{2}
\]

\[
\mathfrak {a} + (\mathfrak {b} + \mathfrak {c}) = (\mathfrak {a} + \mathfrak {b}) + \mathfrak {c}, \quad \mathfrak {a}(\mathfrak {b}\mathfrak {c}) = (\mathfrak {a}\mathfrak {b})\mathfrak {c};\tag{3}
\]

\[
\mathfrak {a} (\mathfrak {b} + \mathfrak {c}) = \mathfrak {a}\mathfrak {b} + \mathfrak {a}\mathfrak {c}.
\]

